\section{Fine-tuning vs. GCD}\label{app_sec:finetuning_vs_gcd}
It is worth considering whether to prefer fine-tuning or GCD to adapt LLMs to structured prediction tasks. Fine-tuning can be applied to either a small model, such as T5~\cite{DBLP:journals/corr/abs-1910-10683}, or directly to a large model, like LLaMA~\cite{touvron2023llama}.
%
There are three crucial factors to take into account:
\begin{enumerate}
    \denselist
    \item availability of training data,
    \item computational cost of fine-tuning, and
    \item performance improvement.
\end{enumerate}
%
Constrained decoding eliminates the need for training data and incurs only the computational cost of running the LM, with small overhead coming from the incremental parser.
%
Fine-tuning a small model is affordable and currently represents the state-of-the-art (SOTA) approach for most structured prediction tasks when training data is available~\cite{decao2021autoregressive, jimenez-gutierrez-etal-2022-thinking}. However, it necessitates a substantial amount of data.
%
Although fine-tuning a large model is expensive, it is highly data-efficient~\cite{brown2020language}.
However, recent advancements in efficient fine-tuning of large models~\cite{hu2021lora} have significantly reduced the computational cost and hardware requirements.
%
We view GCD as a rapid and cost-effective adaptation strategy that enables LMs to produce reliable structured outputs without the need for fine-tuning.
